% Supplementary Methods
\section{Supplementary Methods}\label{sec:supp}

\subsection{ZAP70 as the wet-validation case}\label{sec:supp:zap70}

We use ZAP70, a Syk-family non-receptor tyrosine kinase with long-standing medicinal-chemistry interest but no selective inhibitor in the clinic \cite{auyeung2009zap70}, as the wet-validation anchor for the pipeline. The kinase domain carries a P-loop-adjacent cysteine (Cys346) targetable by an electrophilic warhead. Our anchor is \textbf{Mol1}, an in-house covalent ZAP70 inhibitor with measured pIC$_{50}$ of $6.59$ pairing an isoindoline-acrylamide warhead with a $4$-amino-$1$-isopropyl-imidazole hinge unprecedented in the ZAP70 series; Mol1 trades the canonical two-hydrogen-bond aminopyrimidine hinge for a single-pair anchor whose affinity penalty is offset by covalent residence time \cite{copeland2006residence,bradshaw2015reversible}. We measure (i) per-cohort distributional properties (warhead retention, drug-likeness, predicted potency, covalent-pose feasibility), (ii) recovery and improvement over Mol1 on predicted potency, and (iii) prospective wet-lab activity on selected candidates (in flight at a partner site).

\subsection{Supervised affinity predictor — architecture and evaluation}\label{sec:supp:scorer}

The off-the-shelf affinity predictor invoked as the potency oracle in \S\ref{sec:methods:rl} is, in this work, a FiLM-conditioned $\Delta$pIC$_{50}$ regressor trained on within-assay matched molecular pairs. Given an MMP $(m_a, m_b)$ paired in the same assay, the model computes an edit representation $\delta = \mathrm{fp}(m_b) - \mathrm{fp}(m_a)$ (Morgan fingerprint difference followed by a small MLP encoder), applies the FiLM modulation $\mathrm{FiLM}(h; \delta) = \gamma(\delta) \odot h + \beta(\delta)$ to each endpoint's molecular embedding using \emph{shared} parameters, and produces the predicted $\Delta$pIC$_{50}$ as the difference of two FiLM-modulated reads:
\[
\widehat{\Delta\text{pIC}}_{50}(m_a, m_b) \;=\; g_\phi\bigl(\mathrm{FiLM}(\mathrm{fp}(m_b); \delta)\bigr) - g_\phi\bigl(\mathrm{FiLM}(\mathrm{fp}(m_a); \delta)\bigr).
\]
Training uses $1.7$M within-assay MMPs across $751$ ChEMBL targets. The within-assay constraint is the substantive design choice (assay-level systematic noise cancels in the subtraction); the FiLM architecture is one of several reasonable instantiations. We compared the FiLM-conditioned $\Delta$-predictor against a single-MLP direct-$\Delta$ predictor and a single-molecule absolute regressor on chemistry-disjoint and laboratory-disjoint splits; the comparison is documented separately and is not a contribution of the present paper. The pipeline of \S\ref{sec:methods}--\S\ref{sec:methods:rl} is agnostic to which regressor is used.

\subsection{Generated-cohort filtering}\label{sec:supp:filters}

Sampled molecules from each cohort are screened by a short set of standard medicinal-chemistry filters: validity and warhead retention (SMARTS match for the parent warhead class), Lipinski-style drug-likeness (Lipinski violations $\le 1$, PAINS $< 1$, Brenk reactive alerts $< 2$, SAScore $\le 5$), and Morgan-fingerprint Tanimoto similarity to the anchor ($T_c \ge 0.40$ on the full molecule for the anchor-faithful track; $T_c \ge 0.30$ on the linker-only substructure, warhead and hinge atoms masked, for the scaffold-relaxed track). These filters are standard medicinal-chemistry triage and are not a contribution of the present work.

\subsection{REINVENT4 plumbing details}\label{sec:supp:rl_plumbing}

Two REINVENT4 features that shape the sampling distribution are turned on by default in our TOML configurations. The \texttt{randomize\_smiles} option rewrites the encoder seed at each step as a randomly-rooted canonically-equivalent SMILES of the same molecule, so the policy does not over-fit to a single tokenisation of the anchor. The \texttt{IdenticalMurckoScaffold} diversity filter \cite{bemis1996murcko} keeps a bucketed memory keyed by Bemis--Murcko scaffold and multiplicatively penalises the reward of any candidate whose scaffold has already received $25$ accepted entries (bucket size $25$). The filter caps the effective gradient contribution any one scaffold can make, preventing collapse onto a small set of high-reward chemotypes. Both settings, the DAP hyperparameters of \S\ref{sec:methods:rl}, and the per-component reward scoring thresholds are reproduced verbatim in the TOML configurations released with the code; the staged DAP variants we explored (longer horizons, alternative $\sigma$ schedules, PPO/DPO drop-in replacements) are documented in the same place.

\subsection{Boltz-2 covalent cofold validation}\label{sec:supp:cofold}

Top candidates from each cohort are structurally validated by \boltz~\cite{passaro2025boltz2} co-folding with the ZAP70 kinase domain (UniProt \textsc{P43403}, residues $327$--$606$) under an explicit covalent bond restraint between the warhead $\beta$-carbon (identified via SMARTS) and the Cys346 S$_\gamma$ atom at $1.85$~\r{A}. Because the bond distance is fixed by the restraint, this is a \emph{pose-feasibility} test rather than a free-energy prediction. We record three per-molecule metrics: the \boltz\ confidence (\textsc{iptm}), the S$_\gamma$--C$_\beta$--C$_\alpha$ approach-angle deviation from the canonical Bürgi--Dunitz value (the $95^\circ$--$115^\circ$ thia-Michael transition-state window \cite{burgi1974stereochemistry} serves as the sanity flag), and whether the restraint converged at all.
